\section{Resultados M/M/1}

Los resultados de Python se exportaron a \texttt{results/mm1/experiments/}. El archivo \texttt{mm1\_runs.csv} contiene una fila por corrida, mientras que \texttt{mm1\_summary.csv} y \texttt{mm1\_summary\_long.csv} contienen resumenes por experimento.

\subsection{Cola infinita estable}

Para $\rho=0.25$, $\rho=0.50$ y $\rho=0.75$, la cola infinita tiene regimen estacionario y puede compararse contra las formulas teoricas. La Tabla~\ref{tab:mm1-infinite-steady} muestra que las medias simuladas quedan cerca de los valores teoricos.

% Generado desde results/mm1/experiments/mm1_summary.csv
\begin{table}[htbp]
\centering
\caption{Comparacion teorica y simulada para cola infinita estable.}
\label{tab:mm1-infinite-steady}
\begin{tabular}{llrrrr}
\toprule
$\rho$ & Metrica & Sim. & IC 95\% & Teoria & Error \% \\
\midrule
0.25 & $L$ & 0.3332 & 0.0026 & 0.3333 & -0.04 \\
 & $L_q$ & 0.0829 & 0.0015 & 0.0833 & -0.53 \\
 & $W$ & 0.1329 & 0.0006 & 0.1333 & -0.31 \\
 & $W_q$ & 0.0331 & 0.0005 & 0.0333 & -0.81 \\
 & $\rho_{serv}$ & 0.2503 & 0.0013 & 0.2500 & 0.13 \\
\addlinespace
0.50 & $L$ & 1.0068 & 0.0074 & 1.0000 & 0.68 \\
 & $L_q$ & 0.5052 & 0.0060 & 0.5000 & 1.05 \\
 & $W$ & 0.2009 & 0.0013 & 0.2000 & 0.43 \\
 & $W_q$ & 0.1008 & 0.0011 & 0.1000 & 0.80 \\
 & $\rho_{serv}$ & 0.5015 & 0.0019 & 0.5000 & 0.31 \\
\addlinespace
0.75 & $L$ & 3.0205 & 0.0738 & 3.0000 & 0.68 \\
 & $L_q$ & 2.2705 & 0.0706 & 2.2500 & 0.91 \\
 & $W$ & 0.4025 & 0.0089 & 0.4000 & 0.63 \\
 & $W_q$ & 0.3026 & 0.0087 & 0.3000 & 0.85 \\
 & $\rho_{serv}$ & 0.7500 & 0.0039 & 0.7500 & 0.00 \\
\bottomrule
\end{tabular}
\end{table}


La Tabla~\ref{tab:mm1-queue-probabilities} reporta la probabilidad de observar $q$ clientes esperando en cola. Esta salida proviene del tiempo acumulado por longitud de cola durante la simulacion y se compara contra la distribucion estacionaria correspondiente.

% Generado desde results/mm1/experiments/mm1_queue_probabilities.csv
\begin{table}[htbp]
\centering
\caption{Probabilidad de longitud de cola para cola infinita estable.}
\label{tab:mm1-queue-probabilities}
\begin{tabular}{lccc}
\toprule
$q$ & $\rho=0.25$ & $\rho=0.50$ & $\rho=0.75$ \\
\midrule
$0$ & \shortstack{0.9380\\(0.9375)} & \shortstack{0.7489\\(0.7500)} & \shortstack{0.4372\\(0.4375)} \\
$1$ & \shortstack{0.0466\\(0.0469)} & \shortstack{0.1247\\(0.1250)} & \shortstack{0.1402\\(0.1406)} \\
$2$ & \shortstack{0.0115\\(0.0117)} & \shortstack{0.0628\\(0.0625)} & \shortstack{0.1049\\(0.1055)} \\
$3$ & \shortstack{0.0029\\(0.0029)} & \shortstack{0.0316\\(0.0312)} & \shortstack{0.0786\\(0.0791)} \\
$4$ & \shortstack{0.0008\\(0.0007)} & \shortstack{0.0161\\(0.0156)} & \shortstack{0.0588\\(0.0593)} \\
$5$ & \shortstack{0.0002\\(0.0002)} & \shortstack{0.0079\\(0.0078)} & \shortstack{0.0444\\(0.0445)} \\
\bottomrule
\end{tabular}
\par\smallskip
\begin{minipage}{0.95\linewidth}
\centering\footnotesize
Cada celda muestra simulacion media y, entre parentesis, valor teorico de $P(Q=q)$.
\end{minipage}
\end{table}


\subsection{Cola infinita sin regimen estacionario}

Para $\rho=1.00$ y $\rho=1.25$, la cola infinita no tiene estado estacionario. Las corridas se mantienen para cumplir el barrido experimental, pero sus valores no deben interpretarse como una convergencia teorica.

% Generado desde results/mm1/experiments/mm1_summary.csv
\begin{table}[htbp]
\centering
\caption{Corridas con cola infinita sin regimen estacionario.}
\label{tab:mm1-infinite-unstable}
\begin{tabular}{rrrrrl}
\toprule
$\rho$ & $L$ sim. & $L_q$ sim. & $W$ sim. & $W_q$ sim. & Referencia \\
\midrule
1.00 & 186.3341 & 185.3393 & 18.6264 & 18.5265 & No estacionaria \\
1.25 & 12541.8921 & 12540.8922 & 1003.8583 & 1003.7582 & No estacionaria \\
\bottomrule
\end{tabular}
\end{table}


\subsection{Probabilidad de denegacion en colas finitas}

En colas finitas, todos los valores de $\rho$ tienen referencia estacionaria porque el sistema bloquea arribos cuando alcanza su capacidad total. La Tabla~\ref{tab:mm1-denial-probability} compara la probabilidad de denegacion simulada con la teorica para cada capacidad de cola.

% Generado desde results/mm1/experiments/mm1_summary.csv
\begin{table}[htbp]
\centering
\caption{Probabilidad de denegacion simulada y teorica para colas finitas.}
\label{tab:mm1-denial-probability}
\begin{tabular}{lccccc}
\toprule
$\rho$ & $B=0$ & $B=2$ & $B=5$ & $B=10$ & $B=50$ \\
\midrule
0.25 & \shortstack{0.1994\\(0.2000)} & \shortstack{0.0118\\(0.0118)} & \shortstack{0.0002\\(0.0002)} & \shortstack{0.0000\\(0.0000)} & \shortstack{0.0000\\(0.0000)} \\
0.50 & \shortstack{0.3335\\(0.3333)} & \shortstack{0.0660\\(0.0667)} & \shortstack{0.0081\\(0.0079)} & \shortstack{0.0003\\(0.0002)} & \shortstack{0.0000\\(0.0000)} \\
0.75 & \shortstack{0.4288\\(0.4286)} & \shortstack{0.1549\\(0.1543)} & \shortstack{0.0519\\(0.0513)} & \shortstack{0.0113\\(0.0109)} & \shortstack{0.0000\\(0.0000)} \\
1.00 & \shortstack{0.5009\\(0.5000)} & \shortstack{0.2492\\(0.2500)} & \shortstack{0.1422\\(0.1429)} & \shortstack{0.0836\\(0.0833)} & \shortstack{0.0199\\(0.0192)} \\
1.25 & \shortstack{0.5556\\(0.5556)} & \shortstack{0.3379\\(0.3388)} & \shortstack{0.2533\\(0.2531)} & \shortstack{0.2146\\(0.2148)} & \shortstack{0.1999\\(0.2000)} \\
\bottomrule
\end{tabular}
\par\smallskip
\begin{minipage}{0.95\linewidth}
\centering\footnotesize
Cada celda muestra simulacion media y, entre parentesis, valor teorico.
\end{minipage}
\end{table}


\subsection{Graficos generados}

Los graficos originales se generaron en SVG y luego se copiaron al informe como PDF dentro de \texttt{report/figures/mm1/}. De esta manera se evita depender de soporte SVG adicional al compilar en Overleaf.

\begin{figure}[htbp]
    \centering
    \includegraphics[width=0.82\textwidth]{mm1_infinite_average_number_in_system.pdf}
    \caption{Promedio de clientes en el sistema para cola infinita estable. Solo se incluyen casos con $\rho < 1$, donde existe valor estacionario teorico.}
    \label{fig:mm1-infinite-number-system}
\end{figure}

\begin{figure}[htbp]
    \centering
    \includegraphics[width=0.82\textwidth]{mm1_infinite_average_number_in_queue.pdf}
    \caption{Promedio de clientes en cola para cola infinita estable. La comparacion con teoria se restringe a $\rho < 1$.}
    \label{fig:mm1-infinite-number-queue}
\end{figure}

\begin{figure}[htbp]
    \centering
    \includegraphics[width=0.82\textwidth]{mm1_infinite_average_time_in_system.pdf}
    \caption{Tiempo promedio en el sistema para cola infinita estable.}
    \label{fig:mm1-infinite-time-system}
\end{figure}

\begin{figure}[htbp]
    \centering
    \includegraphics[width=0.82\textwidth]{mm1_infinite_average_time_in_queue.pdf}
    \caption{Tiempo promedio de espera en cola para cola infinita estable.}
    \label{fig:mm1-infinite-time-queue}
\end{figure}

\begin{figure}[htbp]
    \centering
    \includegraphics[width=0.82\textwidth]{mm1_infinite_server_utilization.pdf}
    \caption{Utilizacion del servidor para cola infinita estable.}
    \label{fig:mm1-infinite-utilization}
\end{figure}

\begin{figure}[htbp]
    \centering
    \includegraphics[width=0.82\textwidth]{mm1_denial_probability_by_capacity.pdf}
    \caption{Probabilidad de denegacion de servicio en colas finitas para distintas capacidades de espera $B$.}
    \label{fig:mm1-denial-probability}
\end{figure}

\subsection{Evolucion temporal}

Ademas de las medidas finales, se genero una corrida representativa de cola infinita con $\rho=0.75$ y horizonte de $100$ horas para observar la evolucion del estado del sistema. Esta serie temporal no reemplaza las estimaciones finales de las tablas, que usan $10000$ horas y $10$ replicas, sino que muestra visualmente la dinamica de llegadas, esperas y salidas a lo largo del tiempo.

\begin{figure}[htbp]
    \centering
    \includegraphics[width=0.82\textwidth]{mm1_temporal_clientes_rho_075.pdf}
    \caption{Evolucion temporal de clientes en sistema y clientes en cola para una corrida representativa con $\rho=0.75$.}
    \label{fig:mm1-temporal-clientes}
\end{figure}
